\documentclass{article}

% Language setting
% Replace `english' with e.g. `spanish' to change the document language
\usepackage[czech]{babel}
\usepackage{amsthm,thmtools,xcolor,amsmath,amssymb, mathtools}

% Set page size and margins
% Replace `letterpaper' with `a4paper' for UK/EU standard size
\usepackage[a4paper,top=2cm,bottom=2cm,left=3cm,right=3cm,marginparwidth=1.75cm]{geometry}

% Useful packages
\usepackage{enumerate}
\usepackage{graphicx}
\usepackage{tabularx}
\usepackage[colorlinks=true, allcolors=blue]{hyperref}
\graphicspath{images/}

%\usepackage{tikzit}
%\input{default.tikzstyles}

\title{DÚ Lineární algebra -- Sada 9}
\author{Jan Romanovský}

\begin{document}
\maketitle

\textbf{(9.1)} Převedeme zadání do SLR

$$z_{n+1} = 0,7\cdot z_n + 0,1 \cdot k_n,$$
$$k_{n+1} = - 0,5\cdot z_n + 1,3\cdot k_n,$$

$$\begin{pmatrix}
  z_{n+1} \\ k_{n+1}
\end{pmatrix}=\left(
\begin{array}{cc}
 0,7 & 0,1 \\
 -0,5 & 1,3 \\
\end{array}
\right)\begin{pmatrix}
  z_n \\ k_n
\end{pmatrix},$$

$$\begin{pmatrix}
  z_{n} \\ k_{n}
\end{pmatrix}=\left(
\begin{array}{cc}
 0,7 & 0,1 \\
 -0,5 & 1,3 \\
\end{array}
\right)^n\begin{pmatrix}
  z_0 \\ k_0
\end{pmatrix}.$$

Potom jsou vlastní čísla a vlastní vektory operátoru systému

$$\lambda_1 = 1,2, \, \lambda_2 = 0,8,$$
$$M_{1} = \text{span}\left\{\left(\begin{matrix}
  1 \\ 5
\end{matrix}\right)\right\}, \, M_{2} = \text{span}\left\{\left(\begin{matrix}
  1 \\ 1
\end{matrix}\right)\right\}.$$

Stav systému v čase $n$ je tedy
$$\begin{pmatrix}
  z_n \\ k_n
\end{pmatrix} = \begin{pmatrix}
  1 & 1 \\ 5 & 1
\end{pmatrix}\begin{pmatrix}
  1,2 & 0 \\ 0 & 0,8
\end{pmatrix}^n\begin{pmatrix}
  1 & 1 \\ 5 & 1
\end{pmatrix}^{-1}\begin{pmatrix}
  z_0 \\ k_0
\end{pmatrix}.$$

To se dá také napsat následujícím způsobem

$$\begin{pmatrix} z_n \\ k_n \end{pmatrix} = c_1 \cdot 1,2^n \begin{pmatrix} 1 \\ 5 \end{pmatrix} + c_2 \cdot 0,8^n \begin{pmatrix} 1 \\ 1 \end{pmatrix},$$

kde konstanty vyjádříme pomocí $z_0, k_0$ jako

$$c_1 = \frac{k_0 - z_0}{4}, \quad c_2 = \frac{5z_0 - k_0}{4}.$$

Pro $n\to\infty$ půjde první člen k $c_1 \cdot \infty$ a druhý k nule. Aby populace tedy šla do nekonečna, musí být $c_1$ kladné

$$c_1 > 0 \implies \frac{k_0 - z_0}{4} > 0 \implies k_0 > z_0$$

Počet žab i komárů poroste do nekonečna pro $k_0 > z_0$ a komáři (a po nich i žáby) vymřou pro $k_0 \leq z_0$ (dělící linie těchto dvou případů je $k_0 = z_0$ odpovídá přímce $M_2$).

Limitní poměr komárů ku žabám spočítáme z předchozích vyjádření následující limitou

$$\lim_{n\to\infty} \frac{k_n}{z_n} = \frac{5 \cdot c_1 \cdot 1,2^n}{1\cdot c_1 \cdot 1,2^n} = 5.$$

Pokud počet žab i komárů poroste do nekonečna, poměr komárů ku žabám půjde k $5:1$ (to odpovídá přímce $M_1$).

\end{document}
